\section{E2 - Transformer: Multi-Head Self-Attention và tokenization}

\subsection{Mục tiêu và nguyên tắc hiện thực}
E2 dùng cùng Fashion-MNIST/split với E1. Có hai implementation attention: bản manual là implementation chính; bản \code{nn.MultiheadAttention} chỉ dùng đối chiếu. Không dùng \code{nn.TransformerEncoder} cho bản manual.

\subsection{Multi-Head Self-Attention tự viết}
Với token matrix $X\in\mathbb{R}^{T\times d}$, các projection được học riêng:
\begin{equation}
Q=XW_Q,\quad K=XW_K,\quad V=XW_V,
\end{equation}
\begin{equation}
\mathrm{Attention}(Q,K,V)=\mathrm{softmax}\left(\frac{QK^\top}{\sqrt{d_k}}\right)V.
\end{equation}
Mỗi head được tính độc lập, sau đó concat và đi qua output projection. Báo cáo cần đối chiếu tensor shape của $Q,K,V$, attention weights và output với implementation PyTorch.

\subsection{Ba cách tokenization}
\begin{table}[H]\centering
\caption{Tokenization dùng trong E2.}
\begin{tabularx}{0.94\textwidth}{@{}lXXX@{}}
\toprule
Cách & Số token & Feature/token & Ý nghĩa \\
\midrule
Row & 28 & 28 pixel & Mỗi hàng ảnh là một token \\
Patch $4\times4$ & 49 & 16 pixel & ViT-like non-overlapping patches \\
CNN-stem & 49 & $d_{model}$ & CNN rút feature map $7\times7$, mỗi vị trí spatial là token \\
\bottomrule
\end{tabularx}
\end{table}
\safeimage{e2/e2_tokenization.png}{Minh họa ba cách tokenization trên cùng một ảnh Fashion-MNIST.}

\subsection{Classifier và ma trận thí nghiệm}
Classifier gồm tokenizer, CLS token, hai encoder block nông, LayerNorm và classification head. Để tăng độ sâu so sánh nhưng vẫn bám đề, thực hiện sáu run chính (2 MSA $\times$ 3 tokenizers) và ba run ablation số head 2/4/8 trên tokenizer tốt nhất.

\begin{table}[H]\centering
\caption{Kế hoạch thí nghiệm E2.}
\begin{tabular}{@{}lllrr@{}}
\toprule
Nhóm & MSA & Tokenizer & Heads & Mục đích \\
\midrule
E2-A1--A3 & Manual & Row/Patch/CNN-stem & 4 & So sánh tokenization \\
E2-B1--B3 & PyTorch & Row/Patch/CNN-stem & 4 & Manual vs reference \\
E2-C1--C3 & Manual & Best tokenizer & 2/4/8 & Ablation số head \\
\bottomrule
\end{tabular}
\end{table}

\subsection{Kết quả và attention visualization}
\begin{table}[H]\centering
\caption{Kết quả E2.}
\begin{tabular}{@{}llllrrr@{}}
\toprule
ID & MSA & Tokenizer & Heads & Val Acc & Test Acc & Time/epoch \\
\midrule
\multicolumn{7}{c}{\textcolor{TodoRed}{CHƯA CÓ KẾT QUẢ - SẼ ĐƯỢC NOTEBOOK XUẤT TỪ RUN THẬT}} \\
\bottomrule
\end{tabular}
\end{table}
\safeimage{e2/e2_comparison.png}{So sánh accuracy và thời gian giữa manual/PyTorch MSA và các tokenizer.}
\safeimage{e2/e2_attention_maps.png}{Attention heatmap của một số mẫu Fashion-MNIST.}

\subsection{Thảo luận E2}
\TODO{PHÂN TÍCH: sequence length ảnh hưởng compute ra sao; tokenizer nào tốt; manual và PyTorch có xu hướng tương đồng không; số head giúp/hại thế nào; attention map có hợp lý về mặt thị giác không.}
